\scanpage{31}
% archive sheet number: 30
un des deux éléments de $\varepsilon_T$, i.e. à choisir un \uline{ordre} de cet ensemble, qui nous permet dès lors de distinguer \inserted{un} $L_0$ et un $L_1$, de sorte que
\[
 M=L'_0\oplus L'_1,\tag*{(6)}
\]
et qu'on a un isomorphisme d'identification
\[
 T\xrightarrow{\sim}\unsure{\mathbb G|G}_m,\tag*{(7)}
\]
\inserted{un $\lambda\in G_m(\cdot)\simeq$} $T(\cdot)$ opérant sur $L_0$ par l'homothétie $\lambda$, sur $\unsure{L,|L_1}$ par $\lambda^{-1}$.

On peut écrire, à l'aide de (6)
\[
 H\overset{\mathrm{déf}}{=}\underbrace{\mathrm{Sl}(\deleted{\unsure{2,}}M)\unsure{(\Lambda)}}_{=\operatorname{Aut}(M,e_M)}
 \simeq\left\{\begin{pmatrix}a&b\\c&d\end{pmatrix}\;\middle|\;
 \begin{gathered}
 a,d\in\Lambda,\quad b\in\operatorname{Hom}(L'_1,L'_0)\simeq{L'_1}^{-1}L'_0\\
 c\in\operatorname{Hom}(L'_0,L'_1)\simeq{L'_0}^{-1}L'_1,\quad ad-bc=1
 \end{gathered}\right\},\tag*{(8)}
\]
où \deleted{\unsure{$t$}}
\[
 bc=\deleted{\unsure{b}}cb\in\deleted{\unsure{L_1 L_1|\Lambda^*}}\Lambda^* \tag*{(9)}
\]
désigne le composé de $b$ et $c$, identifié à un élément de $\Lambda$ (homothétie par $\lambda$).

La composition des éléments de $\mathrm{\unsure{SL|Sl}}(\unsure{1|2},M)$ est donné par la régle usuelle. Avec ces notations, l'isomorphisme (7) est donné par
\[
 \Lambda^*\xrightarrow{\sim}T(\Lambda)=\left\{\begin{pmatrix}\lambda&0\\0&\lambda^{-1}\end{pmatrix}\;\middle|\;\lambda\in\Lambda^*\right\}\tag*{(10)}
\]
et on pose $u(\lambda)=\begin{pmatrix}\lambda&0\\0&\lambda^{-1}\end{pmatrix}$ pour $\lambda\in\Lambda^*$.

Par la suite, on s'intéresse au groupe des points \deleted{\unsure{$k$}} de $\mathrm{Sl}(M)$ à valeurs \unsure{ds|de} $\Lambda$, plutôt qu'au schéma en groupes, et on notera par la suite par $T$ au lieu de $T(\Lambda)$, le sous-groupe
